\chapter{深度生成模型}
\begin{quote}\small
\textit{本章摘要。}本章从“学习数据分布而非只学习标签”出发，依次解释变分自编码器、扩散模型及其条件控制，再比较生成质量、覆盖性、可控性和计算代价。ELBO、噪声反演和物理约束把抽象生成目标连接到工业仿真、数据增强和异常样本分析。
\end{quote}

在贯穿案例中，生成模型不能通过制造合成故障来改变真实测试集或故障定义。它只回答数据稀缺时能否补充候选样本、恢复缺失观测或模拟退化轨迹；最终价值仍由未见设备上的真实故障召回、误报、校准和物理约束违反率决定。

本章先说明要学习怎样的联合分布或条件分布，再推导相应的训练目标。接下来检查生成样本是否遗漏某些故障类型、是否符合指定工况，以及采样误差有多大。合成样本可以用于训练，但收益要在保持独立的真实设备测试集上检验。生成的裂纹看起来逼真，也不代表现场出现过这种裂纹；真实故障很少时，模型尤其容易把已有猜测反复复制出来。因此每种方法都要说明哪些性质能够从目标函数得到，哪些仍需观测和验证。

\section{生成建模的目标}
生成模型学习数据分布$p_{\mathrm{data}}(x)$或其条件分布$p(x\mid c)$，因而可以采样、补全、去噪和估计似然。判别模型关注$p(y\mid x)$或决策边界，生成模型还试图描述输入本身的结构。

\paragraph{变分自编码器。}
Variational Autoencoder (VAE)引入潜变量$z$，生成模型为$p_\theta(x,z)=p(z)p_\theta(x\mid z)$\cite{kingma2014vae}。后验$p_\theta(z\mid x)$通常难以计算，以$q_\phi(z\mid x)$近似。由 KL 分解可得
\begin{equation}
 \log p_\theta(x)=\mathcal{L}(x)+\KL(q_\phi(z\mid x)\|p_\theta(z\mid x)),
\end{equation}
其中证据下界（ELBO）为
\begin{equation}
 \mathcal{L}(x)=\E_{q_\phi(z\mid x)}[\log p_\theta(x\mid z)]-\KL(q_\phi(z\mid x)\|p(z)).\label{eq:generative-vae-elbo}
\end{equation}
式\eqref{eq:generative-vae-elbo}右端的第一项重构输入，第二项使潜变量分布接近先验。若$q_\phi(z\mid x)=\mathcal{N}(\mu_\phi(x),\diag(\sigma_\phi^2(x)))$，可用重参数化$z=\mu_\phi(x)+\sigma_\phi(x)\odot\epsilon$、$\epsilon\sim\mathcal{N}(0,I)$进行反向传播。

这个下界可以从 Jensen 不等式直接得到：
\begin{align}
 \log p_\theta(x)
 &=\log\int q_\phi(z\mid x)\frac{p_\theta(x,z)}{q_\phi(z\mid x)}\,dz\\
 &\geq\mathbb E_{q_\phi}[\log p_\theta(x,z)-\log q_\phi(z\mid x)].
\end{align}
将联合分布拆为$p_\theta(x\mid z)p(z)$，就得到重构项减去先验正则项。更完整地说，ELBO的缺口正是后验近似误差：
\begin{align}
\log p_\theta(x)-\mathcal L(x)
&=\KL\bigl(q_\phi(z\mid x)\|p_\theta(z\mid x)\bigr)\geq0.
\end{align}
因此最大化ELBO同时有两个目标：提高观测在生成模型下的解释能力，并让编码器后验逼近真正后验；它并不保证重构误差与工业风险完全一致。若$x\in\mathbb R^D$采用独立高斯解码器$p_\theta(x\mid z)=\mathcal N(\mu_\theta(z),\sigma_x^2I)$，则
\[
-\log p_\theta(x\mid z)=\frac{1}{2\sigma_x^2}\|x-\mu_\theta(z)\|_2^2+\frac D2\log(2\pi\sigma_x^2),
\]
所以常见均方重构项隐含了同方差、条件独立噪声假设。若裂纹像素、高频冲击或少数故障模式更重要，应显式改变观测模型或评价权重，而不能仅增加潜变量维度。

编码器提供近似后验，解码器给出生成机制，两者通过同一个下界联合训练。对高斯后验和标准正态先验，KL 项可解析计算：
\begin{equation}
 \KL(q\|p)=\frac12\sum_{j=1}^{d}(\mu_j^2+\sigma_j^2-\log\sigma_j^2-1).
\end{equation}
若解码器能绕过潜变量解释数据，或 KL 约束过强，近似后验可能接近先验而不再携带输入信息。KL annealing 和 free bits 可以调整训练，但仍需检查潜变量是否保留故障模式，不能只观察总损失。

\section{对抗学习与可逆密度建模}
Generative Adversarial Network (GAN) 的原始对抗目标与可逆密度模型的变量替换目标具有不同的训练信号\cite{goodfellow2014gan,dinh2017realnvp}。本节比较二者如何定义生成分布，而不把样本锐利程度等同于似然质量。
VAE 以可计算的下界换取潜变量推断的便利。若主要目标是采样，可以改用判别器提供训练信号；若需要显式密度，则可以限制变换结构以换取精确的变量替换计算。这两种选择分别引出 GAN 和 Normalizing Flow。

GAN 通过生成器$G$和判别器$D$进行极小极大优化：
\begin{equation}
 \min_G\max_D\;\mathbb{E}_{x\sim p_{data}}\log D(x)+\mathbb{E}_{z\sim p(z)}\log(1-D(G(z))).
\end{equation}
在理想最优判别器下，目标与真实分布和生成分布之间的 Jensen--Shannon 散度（Jensen--Shannon Divergence，JSD）相关。在固定生成分布$p_g$且判别器逐点最优时，对每个$x$令$a=p_{\mathrm{data}}(x)$、$b=p_g(x)$，最大化$a\log D+b\log(1-D)$的一阶导数为$a/D-b/(1-D)$，故
\[
D^*(x)=\frac{p_{\mathrm{data}}(x)}{p_{\mathrm{data}}(x)+p_g(x)}.
\]
代回目标并利用 Jensen--Shannon 散度可得
\[
V(D^*,G)=-\log4+2\operatorname{JSD}(p_{\mathrm{data}}\|p_g),
\]
因此理想博弈的全局最优要求$p_g=p_{\mathrm{data}}$。这一步依赖判别器足够强、分布可比较和优化达到平衡；有限样本或支撑集几乎不相交时，理论散度与实际梯度都可能失真。实际交替训练中，判别器过强可能使原始生成器目标的梯度饱和，模式崩溃则使不同噪声产生相似样本。对工业故障而言，少数典型样本逼真却遗漏罕见模式，仍不足以支持数据增强。

流模型使用可微可逆变换$z=f_\theta(x)$把数据映射到简单先验。这里$x,z\in\R^d$，$J_{f_\theta}(x)$的第$(i,j)$项为$\partial f_{\theta,i}(x)/\partial x_j$，且其行列式非零：
\begin{equation}
 \log p_X(x)=\log p_Z(f_\theta(x))+\log\bigl|\det J_{f_\theta}(x)\bigr|.
\end{equation}
推导来自微分体积变换：若$z=f_\theta(x)$，则$d z=|\det J_f(x)|d x$，而概率守恒给出$p_X(x)dx=p_Z(z)dz$，故$p_X(x)=p_Z(f(x))|\det J_f(x)|$。例如逐维仿射耦合层保持一部分坐标不变，并令另一部分$y_b=x_b\odot\exp s(x_a)+t(x_a)$，其雅可比为三角矩阵，因而$\log|\det J|=\sum_k s_k(x_a)$，同时逆变换可显式计算。可逆性和可计算的 Jacobian 行列式限制了网络设计，也使密度计算无需 VAE 的变分下界。采样时使用逆变换。精确似然只表示对模型分布的精确计算，并不保证该分布符合真实故障机制，更不保证高似然等同于低风险。

表\ref{tab:generative-model-routes}比较的关键是目标与代价是否匹配。需要连续表示时检查 VAE 的压缩损失，需要快速采样时检查 GAN 的模式覆盖，需要密度时检查 Flow 的结构限制。扩散模型则把复杂生成映射分解为多步去噪，后文再推导其训练过程。

\begin{table}[htbp]
\centering
\caption{生成模型路线与工业应用边界}
\label{tab:generative-model-routes}
\scriptsize
\setlength{\tabcolsep}{3pt}
\resizebox{\linewidth}{!}{%
\begin{tabular}{p{2.5cm}p{3.7cm}p{3.7cm}p{4.0cm}}
\hline
模型路线 & 学习对象 & 主要优点 & 工业使用时的关键检查 \\
\hline
VAE & 潜变量分布与重构似然 & 潜空间连续、采样稳定 & 重构是否抹平细裂纹，潜变量是否保留工况 \\
GAN & 生成器与判别器的对抗平衡 & 样本锐利、条件生成直接 & 模式崩溃、训练不稳定、类别覆盖 \\
Flow & 可逆变换与精确似然 & 可计算密度、逆向映射清晰 & 可逆结构限制、显存和表达能力 \\
Denoising Diffusion Probabilistic Model (DDPM)/Denoising Diffusion Implicit Model (DDIM) & 逐步加噪与反向去噪 & 条件控制强、覆盖较好 & 采样步数、异常细节和条件一致性 \\
潜空间扩散 & 压缩空间中的扩散过程 & 高分辨率成本较低 & 编码器压缩损失和生成细节可信度 \\
\hline
\end{tabular}}
\end{table}

\section{扩散过程与条件控制}
\begin{figure}[htbp]
\centering
\begin{tikzpicture}[node distance=0.55cm and 0.45cm]
 \node[idi output, minimum width=1.35cm] (x0) {$x_0$};
 \node[idi process, right=of x0, minimum width=1.35cm] (x1) {$x_1$};
 \node[idi process, right=of x1, minimum width=1.35cm] (x2) {$x_2$};
 \node[idi process, right=of x2, minimum width=1.35cm] (dots) {$\cdots$};
 \node[idi input, right=of dots, minimum width=1.35cm] (xt) {$x_T$};
 \draw[idi arrow] (x0)--node[above,font=\scriptsize]{$+\epsilon$}(x1);
 \draw[idi arrow] (x1)--node[above,font=\scriptsize]{$+\epsilon$}(x2);
 \draw[idi arrow] (x2)--(dots); \draw[idi arrow] (dots)--(xt);
 \draw[idi feedback] (xt.south) .. controls +(0,-0.45) and +(0,-0.45) .. node[below,font=\scriptsize]{学习反向去噪 $\epsilon_\theta$} (x0.south);
\end{tikzpicture}
\caption{扩散模型的时间轴：前向过程逐步破坏结构，反向网络学习在每个噪声尺度恢复结构。}
\label{fig:generative-diffusion-time}
\end{figure}
图\ref{fig:generative-diffusion-time}的实线沿时间增加噪声，反馈箭头则概括学习到的反向恢复路径；训练需覆盖多个噪声尺度，而不是只拟合从末端噪声到样本的一次映射。
扩散模型通过逐步加噪和反向去噪构造生成过程\cite{ho2020ddpm}。前向过程逐步加入高斯噪声：
\begin{equation}
 q(x_t\mid x_{t-1})=\mathcal{N}(\sqrt{1-\beta_t}x_{t-1},\beta_tI).
\end{equation}
这里$t=1,\ldots,T$为扩散步，$0<\beta_t<1$为噪声日程，$x_0\sim p_{\mathrm{data}}$。记$\alpha_t=1-\beta_t$、$\bar\alpha_t=\prod_{s=1}^{t}\alpha_s$，则可直接采样
\begin{equation}
 x_t=\sqrt{\bar\alpha_t}x_0+\sqrt{1-\bar\alpha_t}\epsilon,
 \qquad \epsilon\sim\mathcal{N}(0,I).\label{eq:generative-direct-noising}
\end{equation}
反向网络$\epsilon_\theta(x_t,t)$预测噪声，常用简化目标
\begin{equation}
 \mathcal{L}_{\mathrm{simple}}=\E_{t,x_0,\epsilon}\left[\|\epsilon-\epsilon_\theta(x_t,t)\|_2^2\right].
\end{equation}
此处独立采样$t\sim\operatorname{Uniform}\{1,\ldots,T\}$、$x_0\sim p_{\mathrm{data}}$和$\epsilon\sim\mathcal N(0,I)$，再由式\eqref{eq:generative-direct-noising}构造$x_t$。由条件期望的最小均方性质，若网络容量和优化均理想，则
\[
\epsilon_\theta^*(x_t,t)=\mathbb E[\epsilon\mid x_t,t].
\]
又因$x_t=\sqrt{\bar\alpha_t}x_0+\sqrt{1-\bar\alpha_t}\epsilon$，可由噪声预测得到干净样本估计
\[
\widehat x_0=\frac{x_t-\sqrt{1-\bar\alpha_t}\,\epsilon_\theta(x_t,t)}{\sqrt{\bar\alpha_t}}.
\]
当$\bar\alpha_t$很小时，这个反演会放大预测误差，说明末端高噪声阶段对模型校准和采样数值稳定性尤其敏感。

采样时从高斯噪声开始，按反向时间步逐渐去噪。直观上，模型学习的不是一次性画出样本，而是学习一系列局部的“如何去掉一点噪声”。

\paragraph{扩散模型的后验。}
由于前向过程是高斯马尔可夫链，给定$x_0,x_t$的后验仍为高斯：
\begin{equation}
 q(x_{t-1}\mid x_t,x_0)=\mathcal{N}(\tilde\mu_t(x_t,x_0),\tilde\beta_tI).
\end{equation}
对$t=2,\ldots,T$，其参数为
\begin{align}
 \tilde\mu_t(x_t,x_0)&=\frac{\sqrt{\bar\alpha_{t-1}}\beta_t}{1-\bar\alpha_t}x_0
 +\frac{\sqrt{\alpha_t}(1-\bar\alpha_{t-1})}{1-\bar\alpha_t}x_t,\label{eq:generative-posterior-mean}\\
 \tilde\beta_t&=\frac{1-\bar\alpha_{t-1}}{1-\bar\alpha_t}\beta_t.
\end{align}
其推导只需把两个高斯因子相乘。由
$q(x_t\mid x_{t-1})\propto\exp[-\|x_t-\sqrt{\alpha_t}x_{t-1}\|^2/(2\beta_t)]$
以及
$q(x_{t-1}\mid x_0)=\mathcal N(\sqrt{\bar\alpha_{t-1}}x_0,(1-\bar\alpha_{t-1})I)$，收集关于$x_{t-1}$的二次项，精度为
$\beta_t^{-1}+(1-\bar\alpha_{t-1})^{-1}$，所以方差是其倒数$\tilde\beta_t$；再收集一次项并乘以该方差，就得到式\eqref{eq:generative-posterior-mean}的均值。这说明后验同时利用当前噪声状态$x_t$和原始状态$x_0$。采样时$x_0$未知，网络必须先由$x_t$估计它或等价地估计噪声，不能把该后验公式误当作可直接使用的生成器。
约定$\bar\alpha_0=1$，则$t=1$时方差为零，条件后验退化到已知$x_0$，不作为非退化高斯参与上述 KL 计算。
对$t\geq2$，扩散模型的变分目标包含上述后验与反向模型$p_\theta(x_{t-1}\mid x_t)$之间的 KL 项。固定反向方差并采用噪声预测参数化后，这些项可写为带时间步权重的噪声预测误差；常用简化目标改变了这些权重，不能与完整变分目标直接等同。采样时无法获得$x_0$，因此反向模型必须仅根据当前带噪状态和时间步完成预测。噪声预测、$x_0$预测和速度预测是不同参数化，改变梯度尺度与采样表现，但共享逐步去噪框架。

\paragraph{条件扩散与引导。}
若有条件$c$，网络写为$\epsilon_\theta(x_t,t,c)$。分类器引导使用带噪样本分类器的梯度调整生成方向；无分类器引导通过训练时随机丢弃条件，使网络同时学习条件与无条件预测，再组合为
\begin{equation}
 \hat\epsilon=\epsilon_\theta(x_t,t,\varnothing)
 +w[\epsilon_\theta(x_t,t,c)-\epsilon_\theta(x_t,t,\varnothing)].
\end{equation}
$w=0$对应无条件预测，$w=1$对应原始条件预测，$w>1$进一步加强条件方向。更强引导可能损害多样性或产生伪影，应在固定采样预算下比较条件满足率和缺陷覆盖。ControlNet 等条件分支可接收边缘、深度或几何信息，但条件输入并不自动保证生成结果满足物理约束。

\paragraph{Score Matching。}
以下$p_t(x_t)=\int q(x_t\mid x_0)p_{\mathrm{data}}(x_0)\,dx_0$表示前向带噪边缘密度，不是反向模型的转移密度。
扩散模型也可从各噪声时刻的 score $\nabla_{x_t}\log p_t(x_t)$理解。去噪 score matching 估计带噪分布的对数密度梯度，而不直接计算归一化常数。高斯加噪下，噪声预测对应的 score 估计为$-\epsilon_\theta(x_t,t)/\sqrt{1-\bar\alpha_t}$。这一估计需与噪声日程共同构成反向随机过程，或对应的概率流常微分方程，才能用于分布采样；单纯沿密度梯度上升会趋向局部高密度区域，并不保证恢复目标分布的模式比例。这个观点连接了扩散模型、能量模型和随机微分方程。

\section{采样、约束与工业可信度}
DDIM 在与 DDPM 相关的训练目标下改变采样轨迹，潜空间扩散通过压缩表示降低高分辨率生成成本\cite{song2021ddim,rombach2022latent}。无分类器引导与 ControlNet 分别提供条件强度组合和结构条件分支\cite{ho2022cfg,zhang2023controlnet}，两者均不能自动保证物理可行性。
有了可采样模型，还不能直接把生成样本当成维护数据。只用正常数据训练的模型可以模仿正常模式，却没有凭据知道故障长什么样。若要按故障标签或严重度生成，应有真实故障样本、经过验证的仿真机制，或能够核查的领域约束。不同用途也要分别验收：补足少数类别要检查模式是否被覆盖，模拟退化要检查时间顺序和工况，指定位置的缺陷图像要检查位置与标注是否一致。

\paragraph{DDIM 与采样加速。}
DDPM 使用随机反向过程，需要较多时间步。DDIM 可以构造非马尔可夫的确定性或弱随机采样轨迹，在保留训练目标的同时减少采样步数。步数减少会降低延迟，却可能损害细节和条件一致性，因此应绘制质量--延迟曲线，而不是只报告单一采样配置。

\paragraph{潜空间扩散。}
直接在高维像素或原始波形空间扩散成本很高。潜空间扩散先用编码器把输入压缩到$z$，在潜空间去噪，再由解码器重构。压缩比例过大可能删除细微裂纹、冲击脉冲或高频异常，因此工业应用必须比较原空间频谱、边缘和峰值统计，而不能只看重构图像的主观质量。

\begin{figure}[htbp]
\centering
\begin{tikzpicture}[>=Latex, node distance=1.15cm, every node/.style={font=\small}]
 \node[draw, rounded corners, fill=blue!10, minimum width=2.3cm] (x) {$x_0$：真实样本};
 \node[draw, rounded corners, fill=orange!12, right=of x, minimum width=2.5cm] (noise) {\shortstack{前向加噪\\$q(x_t|x_0)$}};
 \node[draw, rounded corners, fill=gray!15, right=of noise, minimum width=2.3cm] (xt) {$x_T$：近似噪声};
 \draw[->, thick] (x)--(noise); \draw[->, thick] (noise)--(xt);
 \draw[->, thick, dashed] (xt.north) .. controls +(0,0.65) and +(0,0.65) .. node[above,font=\scriptsize]{学习反向去噪} (x.north);
 \node[draw, rounded corners, fill=green!12, below=1.1cm of noise, minimum width=3cm] (net) {$\epsilon_\theta(x_t,t,c)$};
 \draw[->] (noise.south)--(net.north);
\end{tikzpicture}
\caption{扩散模型把生成拆为前向加噪和学习反向去噪。}
\label{fig:generative-denoising-network}
\end{figure}
图\ref{fig:generative-denoising-network}把预设加噪机制与可训练的噪声预测网络分开；网络读取带噪状态、时间和条件，采样时再由这些预测逐步恢复样本。

\paragraph{物理约束生成。}
工业生成样本必须满足守恒、边界、单调性、频率范围或设备状态转移约束。可以在损失中加入残差惩罚
\begin{equation}
 \mathcal{L}=\mathcal{L}_{gen}+\lambda_{phys}\mathbb{E}[\|r(x,c)\|_2^2],
\end{equation}
也可以在采样后投影到可行域。惩罚过弱会产生不可能样本，过强则可能牺牲数据多样性；因此应分别报告真实性、多样性和约束违反率。

\paragraph{记忆、隐私与生成数据。}
生成模型可能记住训练集中的罕见记录，尤其在数据量小、重复度高或模型过拟合时。近邻距离、成员推断和重复片段检测可用于筛查记忆风险。差分隐私训练可以降低单个样本的影响，但通常带来质量--隐私权衡。生成数据用于发布或跨机构共享前，应进行去标识、相似样本审计和下游泄漏测试。

\paragraph{生成模型的实验流程。}
完整实验应包含训练数据分布、采样配置、随机种子、生成数量、重复检测、专家评价、统计检验和下游任务。对故障生成尤其应使用未参与生成训练的真实故障测试集；否则增强后的模型可能只是在适应合成数据的风格。最有说服力的结论是：在不同随机种子、不同设备和不同故障类型上，真实测试性能都获得稳定改善。

评价应先检查边缘分布、相关性、频谱与缺陷几何，再检查物理约束和条件一致性，最后比较真实数据上的下游结果。Fréchet Inception Distance (FID) 或主观逼真度只覆盖部分属性。对照应至少包括仅真实数据与真实加合成数据，并匹配训练预算；重采样基线可以区分数量效应，打乱合成标签则可作为诊断对照，不能替代有效基线。

新增样本可能复制设备身份和批次标记，也可能改变故障先验比例。因此测试集须保持真实、独立且未参与生成器训练与选择，报告召回、误报和校准，而不只看训练损失。若训练指标上升而真实测试性能下降，应检查覆盖不足、生成伪影和先验变化，不能继续用更多合成数据掩盖问题。

\paragraph{本章小结。}
VAE 借助潜空间生成，GAN 通过对抗判别训练，flow 使用可逆变换计算似然，diffusion 逐步去噪得到样本。它们的计算方式不同，生成结果都仍需检查。用于工业任务时，既看样本是否逼真、是否覆盖需要的故障，也检查是否违背物理规律、是否泄露训练记录，最后用独立的真实数据检验下游收益。